% 原书第 22—25 页；第 2.3.3—2.3.5 小节及第 2.4 节。
\subsection{建立可识别性}
\label{subsec:ch02-identifiability}

我们已经作出的假设，是否足以将感兴趣的因果参数，表示为观测数据概率分布 $P_0$ 的一个参数？我们希望，对某个参数映射 $\Psi$，能够将 $\Psi^F(P_{U,X,0})$ 写为 $\Psi(P_0)$。这里提醒读者，SCM 也规定了：$(U,X)$ 的真实分布 $P_{U,X,0}$ 如何导出观测数据结构 $O$ 的分布 $P_0$。由于 $(U,X)$ 的真实概率分布可以是 SCM $\mathcal M^F$ 中的任意元素，而每一个这样的选择 $P_{U,X}$ 都会导出 $O$ 的一个概率分布 $P(P_{U,X})$，因此，我们需要证明，对于所有 $P_{U,X}\in\mathcal M^F$，都有 $\Psi^F(P_{U,X})=\Psi(P(P_{U,X}))$。

这一步需要对 $U$ 的分布，或者有时还要对确定性函数 $f$，作出可能需要的额外假设，使我们能够从观测数据分布中识别目标参数。因此，对于底层数据 $(U,X)$ 的每个概率分布，只要它满足 SCM 以及这些可能施加于 $P_U$ 的额外假设，就存在某个 $\Psi$，使 $\Psi^F(P_{U,X})=\Psi(P(P_{U,X}))$。$O$ 由 $(U,X)$ 的分布所导出，例如 $O=X$ 或 $O\subset X$，且 $P=P(P_{X,U})$；其中，$P(P_{U,X})$ 是 $P_{U,X}$ 所导出的 $O$ 的分布。

将这样得到的完整数据 SCM 记为 $\mathcal M^{F*}\subset\mathcal M^F$，是为了明确：我们可能加入了一些纯粹由可识别性问题驱动的额外假设，它们未必反映现实。确切地说，$\mathcal M^F$ 是在基于真实知识的假设下建立的完整数据 SCM；$\mathcal M^{F*}$ 则是在目标参数的可识别性所需的额外因果假设下建立的完整数据 SCM。现在，对于每个 $P_{U,X}\in\mathcal M^{F*}$，都有 $\Psi^F(P_{U,X})=\Psi(P)$，其中 $P=P(P_{U,X})$ 是 $P_{U,X}$ 所导出的 $O$ 的分布；与此相对，$P_0$ 是真实分布 $P_{U,X,0}$ 所导出的 $O$ 的真实分布。

已有一些定理，可以帮助我们建立所需的可识别性结果。对于作用于 $A$ 节点的某个特定干预 $d$，以及一个给定的实现值 $u$，SCM 会确定地生成 $L_{1:2}$ 的相应取值。将所得实现值记为 $L_d(u)$；注意，$L_d(u)$ 由 $f$ 和 $u$ 所确定。如果 $O=X$，并且 $U$ 的分布满足：在每个时点 $t$，给定 $\operatorname{Pa}(A_t)$ 时，$A_t$ 与 $L_d$ 独立，那么，g 公式（g-formula）就能用 $O$ 的分布来表示 $L_d$ 的分布：
% 原书条件式在 Pa_d(L_t) 后多印一个右括号；按印刷原样保留。
\[
 P(L_d=l)=\prod_{t=1}^{T}
 P(L_t=l_t\mid\operatorname{Pa}_d(L_t))=\operatorname{Pa}_d(l_t)),
\]
其中，$\operatorname{Pa}_d(L_t)$ 是 $L_t$ 的父节点，而这些父节点中的干预节点，由干预 $d$ 确定地设定取值。对于 $U$ 的某种特定独立结构，可以通过在相应的因果图上检验后门路径准则（backdoor path criterion），来建立这一所谓的序贯随机化假设（sequential randomization assumption）；这里的因果图由 SCM 及 $U$ 的这一独立结构共同导出。后门路径准则指出：对于每个 $A_t$，从 $A_t$ 出发、通向在 $A_t$ 之后实现的某个 $L_t$ 节点的每条后门路径，都被其他某个 $L_t$ 节点阻断。

通过这种方式，我们可能得到关于 $U$ 分布的若干独立结构，它们都能给出所需的可识别性结果。也就是说，能够给出所需可识别性的 $U$ 的模型，可以表示为若干模型的并集，而其中每个模型都假定 $U$ 具有某种特定的独立结构。序贯随机化假设也称为无未测量混杂假设。我们将混杂因素（confounder）定义为 $X$ 中这样一些变量：为了建立目标参数的可识别性，我们必须在 $O$ 中观测它们。需要注意的是，$X$ 的不同这类子集，都可能给出所需的可识别性结果。回到 IVF 研究及式 (2.1) 中的结构方程模型，若干独立结构的并集能够满足后门路径准则，从而识别我们的因果目标参数。图 2.3 中的独立结构不满足后门路径准则。

\subsection{选定统计模型与目标参数}
\label{subsec:ch02-statistical-model-target}

可识别性结果为我们提供了一个纯粹的统计目标参数 $\Psi(P_0)$，它定义在 $O$ 的分布 $P_0$ 上。完整数据模型 $\mathcal M^{F*}$ 导出一个观测数据统计模型 $\mathcal M=\{P(P_{X,U}):P_{X,U}\in\mathcal M^{F*}\}$，用于描述 $O$ 的分布 $P_0=P(P_{U,X,0})$。这样就定义了目标参数 $\Psi\colon\mathcal M\to\mathbb R^d$。对于 $\mathcal M^F$ 和 $\mathcal M^{F*}$，$O$ 分布的观测数据统计模型可能相同。如果不同，我们可以考虑将 $\Psi$ 扩展到 $\mathcal M^F$ 所导出的、更大的观测数据统计模型，例如一个允许所有概率分布的完全非参数模型。这样，如果更受限的 SCM 成立，我们的目标参数仍然对应所要估计的目标参数；同时，也允许数据与那个建立在额外可疑假设上的、更受限的 SCM 相矛盾。

对于 IVF 研究，相应的统计参数 $\Psi(P_0)$ 为
% 与原书第 21 页及第 1 章相同的无下标求和和 L_1=1，均照录。
\[
 \begin{aligned}
 \Psi^F(P_{U,X,0})&=P(Y_{(1,1)}=1)\\
 &=E_{P_0}\Bigl(\sum P_0(Y=1\mid A_{0:1}=1,L_{0:1}=l_{0:1})\\
 &\hspace{3.2em}{}\times P_0(L_1=1\mid A_0=1,L_0=l_0)\Bigr)
 \equiv\Psi(P_0).
 \end{aligned}
\]
为了得到这一可识别性结果，将因果效应表示为 $O$ 的分布 $P_0$ 的一个参数，我们需要作出序贯随机化假设。这一假设可能已经包含在最初的 SCM $\mathcal M^F$ 中；但如果已知存在未测量的混杂因素，那么，模型 $\mathcal M^{F*}$ 就会通过强制施加一个我们认为不正确的随机化假设，变得更加受限。

另一个必需的假设是序贯正值性（sequential positivity）。在 IVF 研究中，这意味着：给定协变量历史，在两个时点中的每个时点，接受 IVF 治疗的概率都非零。没有这一假设，$\Psi(P_0)$ 中涉及 $L_{1:2}$ 的概率就没有良好定义。不过，“正值性假设”是一个更一般的名称，指使目标参数 $\Psi(P_0)$ 定义良好所必需的条件；它往往要求删失机制或处理机制具有某种特定的支撑集（support）。

\subsection{目标参数的解释}
\label{subsec:ch02-interpretation}

我们可能没有足以支持 SCM 中因果假设的知识，也可能不愿依赖这些额外假设。只要假设观测变量 $L_{0:1}$、$A_{0:1}$ 和 $Y$ 的时间顺序正确，即
\[
 L_0\to A_0\to L_1\to A_1\to Y,
\]
我们的目标参数仍然表示一个有意义、定义良好的效应，并且可以解释为变量重要性度量（variable importance measure）。

\begin{sourcebox}
观测数据参数 $\Psi(P_0)$ 可以有两种可能不同的解释：
\begin{enumerate}
 \item 在 $P_0\in\mathcal M$ 的基础上，加入那些真正可靠且已知成立的额外非统计假设（例如 $\mathcal M^F$），来解释 $\Psi(P_0)$。这可能涉及：在一个现实的因果模型 $\mathcal M^F$ 下，为 $\Psi(P_0)$ 相对于所需目标因果效应 $\Psi^F(P_{U,X,0})$ 的偏离程度给出界，而这一因果模型尚不足以识别该因果效应。
 \item 在更受限的 SCM $\mathcal M^{F*}$ 下，将其解释为真正的因果参数 $\Psi^F(P_{U,X})=\Psi(P_0)$。此时，也就纳入了从 $O$ 的概率分布 $P_0$ 识别所需因果效应所必需的全部因果假设。
\end{enumerate}
\end{sourcebox}

\section{注记与延伸阅读}
\label{sec:ch02-notes}

关于 SCM 的深入介绍，可参阅 Pearl（2009a）。本章以《靶向学习》（2011）第 2 章和第 24 章为基础，并依赖其中的内容。部分内容经许可再次采用了该书第 2 章的材料。在序贯随机化假设下，从观测数据分布识别反事实分布的 g 公式，最早发表于 Robins（1986）。
